\documentclass[11pt,a4paper]{article}
\usepackage[utf8]{inputenc}
\usepackage[indonesian]{babel}
\usepackage{amsmath,amsfonts,amssymb,amsthm}
\usepackage{graphicx}
\usepackage{geometry}
\usepackage{fancyhdr}
\usepackage{tcolorbox}
\usepackage{booktabs}
\usepackage{tabularx}
\usepackage{tikz}
\usetikzlibrary{positioning,shapes.geometric,arrows.meta,calc}
\usepackage{circuitikz}
\usepackage{enumitem}
\usepackage{listings}
\usepackage{xcolor}
\usepackage{hyperref}

\geometry{top=2.5cm, bottom=2.5cm, left=2.5cm, right=2.5cm, headheight=15pt}

\pagestyle{fancy}
\fancyhf{}
\fancyhead[L]{\small\textbf{Persamaan Diferensial 2026} -- Teknik Elektro UNIB}
\fancyhead[R]{\small Modul 3: Aplikasi Transien RC, RL \& Termal Trafo}
\fancyfoot[C]{\thepage}

% Pengaturan kode Julia
\lstdefinelanguage{Julia}%
  {morekeywords={abstract,break,case,catch,const,continue,do,else,elseif,%
      end,export,false,for,function,global,if,import,%
      let,local,macro,module,mutable,primitive,quote,return,struct,%
      true,try,type,using,var,while},%
   sensitive=true,%
   morecomment=[l]{\#},%
   morecomment=[n]{\#=}{=\#},%
   morestring=[b]',%
   morestring=[b]"%
  }[keywords,comments,strings]

\lstdefinestyle{juliastyle}{
    language=Julia,
    basicstyle=\footnotesize\ttfamily,
    keywordstyle=\color{blue!80!black}\bfseries,
    commentstyle=\color{green!50!black}\itshape,
    stringstyle=\color{red!80!black},
    numbers=left,
    numberstyle=\tiny\color{gray},
    breaklines=true,
    showstringspaces=false,
    frame=single,
    backgroundcolor=\color{gray!6},
    captionpos=b,
    xleftmargin=10pt,
    tabsize=4
}

\definecolor{headercolor}{RGB}{0,51,102}
\definecolor{accentcolor}{RGB}{0,102,204}
\definecolor{shadecolor}{RGB}{245,248,253}

\hypersetup{
    colorlinks=true,
    linkcolor=headercolor,
    citecolor=accentcolor,
    urlcolor=blue
}

\theoremstyle{definition}
\newtheorem{definition}{Definisi}[section]
\newtheorem{example}{Contoh Soal}[section]
\newtheorem{theorem}{Teorema}[section]

\title{\textbf{\Large MODUL AJAR KULIAH MINGGU 3}\\[0.3cm]
\textbf{\LARGE APLIKASI PERSAMAAN DIFERENSIAL BIASA ORDE 1 DALAM REKAYASA ELEKTRO: DINAMIKA TRANSIEN RC-RL, LONJAKAN TEGANGAN INDUKTIF (\textit{INDUCTIVE KICKBACK}) \& PEMODELAN TERMAL TRANSFORMATOR DAYA DENGAN HUKUM PENDINGINAN NEWTON}}
\author{\textbf{Ir. Novalio Daratha, S.T., M.Sc., Ph.D.} \& \textbf{Muhammad Arfan, S.T., M.T.} \\ 
Jurusan Teknik Elektro -- Fakultas Teknik \\ 
Universitas Bengkulu}
\date{Semester Genap 2026}

\begin{document}

\maketitle

\begin{center}
  \vspace{-0.25cm}
  \begin{tcolorbox}[colback=red!4!white,colframe=red!75!black,boxrule=0.5pt,arc=2pt,left=6pt,right=6pt,top=3pt,bottom=3pt,width=0.98\textwidth]
    \centering\footnotesize
    \textbf{Video Perkuliahan Resmi (YouTube):} {\color{red!80!black}\textbf{\url{https://youtu.be/_JZ7kBkgxlk}}} \quad $\bullet$ \quad \textbf{Playlist:} \href{https://www.youtube.com/playlist?list=PLS5oOZWXZeTw}{\color{blue!80!black}\textbf{bit.ly/pd-unib-playlist}} \quad $\bullet$ \quad \textbf{Portal:} \url{https://pd.ndaratha.my.id}
  \end{tcolorbox}
  \vspace{-0.15cm}
\end{center}

\begin{tcolorbox}[colback=blue!5!white,colframe=headercolor,title=\textbf{Capaian Pembelajaran Khusus (Sub-CPMK 2)}]
Setelah menyelesaikan pembelajaran pada Modul Ajar Minggu 3 ini, mahasiswa diharapkan mampu:
\begin{enumerate}[leftmargin=*]
    \item \textbf{C1 (Mengingat):} Menyatakan bentuk matematis Masalah Nilai Awal (IVP) untuk rangkaian transien RL, RC, dan model kesetimbangan energi termal berbasis Hukum Pendinginan Newton.
    \item \textbf{C2 (Memahami):} Menjelaskan prinsip kontinuitas arus induktor ($i_L(0^-)=i_L(0^+)$), fenomena lonjakan tegangan induktif (\textit{inductive kickback}) saat pemutusan arus mendadak, serta konsep konstanta waktu listrik ($\tau = L/R$, $\tau = RC$) dan konstanta waktu termal transformator ($\tau_{\text{th}}$).
    \item \textbf{C3 (Menerapkan):} Menurunkan solusi analitik analitis langkah demi langkah untuk respon alami (\textit{natural response}) dan respon paksa (\textit{step response}) pada rangkaian RL, RC multi-interval, serta profil suhu transien transformator daya akibat pembebanan dinamis.
    \item \textbf{C4 (Menganalisis):} Menganalisis kurva pengisian dan pengosongan energi elektromagnetik, transien sakelar daya semikonduktor, serta dampak beban lebih darurat terhadap laju penuaan relatif insulasi minyak dan kertas trafo berdasarkan model Arrhenius.
    \item \textbf{C5 (Mengevaluasi):} Mengevaluasi efektivitas sirkuit proteksi dioda \textit{freewheeling} dalam membatasi busur api (\textit{electric arcing}) pada kontak relai dan menentukan ambang batas aman pembebanan transformator gardu induk PLN sesuai standar IEEE Std C57.91.
    \item \textbf{C6 (Menciptakan/Komputasi):} Mengembangkan skrip simulasi numerik terbuka berbasis \textbf{Julia} (\texttt{DifferentialEquations.jl} \& \texttt{Plots.jl}) untuk memodelkan respon transien elektrik dan kurva kenaikan suhu termal transformator daya secara simultan.
\end{enumerate}
\end{tcolorbox}

\vspace{0.3cm}
\tableofcontents
\vspace{0.5cm}

\newpage

\section{Peta Konsep \& Posisi Modul dalam Kurikulum}

Dalam kurikulum 16 minggu mata kuliah Persamaan Diferensial di Jurusan Teknik Elektro Universitas Bengkulu, Minggu 3 menandai jembatan krusial dari fondasi matematis murni (Minggu 1 dan 2) menuju penerapan rekayasa tingkat lanjut. Pada tahap ini, seluruh teknik pengintegrasian PDB linier orde 1 (separabel, eksak, dan faktor pengintegrasi) diterapkan secara simultan untuk membedah dua pilar utama teknik elektro modern:
\begin{itemize}[leftmargin=*]
    \item \textbf{Transien Elektromagnetik Rangkaian Terpusat:} Penyimpanan dan pelepasan energi pada elemen reaktif $L$ dan $C$, fenomena \textit{switching}, dan mitigasi lonjakan tegangan balik destruktif (\textit{inductive kickback}).
    \item \textbf{Dinamika Termal Sistem Tenaga Listrik:} Disipasi rugi-rugi daya Joule ($I^2 R$) dan histeresis/arus eddy pada inti besi transformator gardu induk yang memicu pemanasan minyak isolasi, dimodelkan secara elegan melalui analogi rangkaian termal berbasis Hukum Pendinginan Newton.
\end{itemize}

Posisi strategis Modul Minggu 3 di dalam keseluruhan kurikulum disajikan pada Gambar~\ref{fig:peta_jalan_pd_w3}.

\begin{figure}[htbp]
\centering
\resizebox{\linewidth}{!}{%
\begin{tikzpicture}[
    node distance=0.85cm and 0.45cm,
    block/.style={rectangle, draw=headercolor, fill=blue!8, thick, rounded corners=3pt, align=center, text width=3.35cm, minimum height=1.05cm, font=\scriptsize},
    doneblock/.style={rectangle, draw=green!60!black, fill=green!10, thick, rounded corners=3pt, align=center, text width=3.35cm, minimum height=1.05cm, font=\scriptsize},
    highlight/.style={rectangle, draw=red!80!black, fill=red!15, very thick, rounded corners=3pt, align=center, text width=3.35cm, minimum height=1.05cm, font=\scriptsize\bfseries},
    midblock/.style={rectangle, draw=orange!80!black, fill=orange!15, thick, rounded corners=3pt, align=center, text width=3.35cm, minimum height=1.05cm, font=\scriptsize\bfseries},
    arrow/.style={->, >=stealth, line width=1.1pt, headercolor}
]
    % Paruh 1: PDB & Rangkaian Listrik
    \node[doneblock] (w1) {Minggu 1: Klasifikasi PD, IVP\\\& Pemodelan Fisis RL};
    \node[doneblock, right=of w1] (w2) {Minggu 2: PDB Orde 1\\(Separabel \& Faktor Integrasi)};
    \node[highlight, right=of w2] (w3) {Minggu 3: Aplikasi Rekayasa\\(Transien RC, RL, Termal)};
    \node[block, right=of w3] (w4) {Minggu 4: PDB Orde 2 Homogen\\(Wronskian \& Tangki LC)};
    
    \node[block, below=of w4] (w5) {Minggu 5: PDB Orde 2 Non-Homogen\\\& Transien RLC Seri AC};
    \node[block, left=of w5] (w6) {Minggu 6: Transformasi Laplace\\Dasar \& Teorema Geser};
    \node[block, left=of w6] (w7) {Minggu 7: Invers Laplace \&\\Solusi PDB Domain $s$};
    \node[midblock, left=of w7] (w8) {Minggu 8: Evaluasi Capaian\\Ujian Tengah Semester (UTS)};
    
    % Paruh 2: PDP & Medan Elektromagnetika
    \node[block, below=of w8] (w9) {Minggu 9: Pengantar PDP \&\\Analisis Deret Fourier};
    \node[block, right=of w9] (w10) {Minggu 10: Solusi PDP Metode\\Pemisahan Variabel};
    \node[block, right=of w10] (w11) {Minggu 11: Persamaan Gelombang\\1D (Telegrafer Transmisi)};
    \node[block, right=of w11] (w12) {Minggu 12: Persamaan Panas 1D\\\& Manajemen Termal Kabel};
    
    \node[block, below=of w12] (w13) {Minggu 13: Persamaan Laplace 2D\\Potensial Elektrostatik};
    \node[block, left=of w13] (w14) {Minggu 14: Fungsi Bessel \&\\Efek Kulit (\textit{Skin Effect})};
    \node[block, left=of w14] (w15) {Minggu 15: 4 Persamaan Maxwell\\Diferensial \& Sintesis PDP};
    \node[midblock, left=of w15] (w16) {Minggu 16: Evaluasi Akhir\\Ujian Akhir Semester (UAS)};
    
    \draw[arrow] (w1) -- (w2);
    \draw[arrow] (w2) -- (w3);
    \draw[arrow] (w3) -- (w4);
    \draw[arrow] (w4) -- (w5);
    \draw[arrow] (w5) -- (w6);
    \draw[arrow] (w6) -- (w7);
    \draw[arrow] (w7) -- (w8);
    \draw[arrow] (w8) -- (w9);
    \draw[arrow] (w9) -- (w10);
    \draw[arrow] (w10) -- (w11);
    \draw[arrow] (w11) -- (w12);
    \draw[arrow] (w12) -- (w13);
    \draw[arrow] (w13) -- (w14);
    \draw[arrow] (w14) -- (w15);
    \draw[arrow] (w15) -- (w16);
\end{tikzpicture}%
}
\caption{Peta jalan kurikulum perkuliahan dengan penekanan pada Modul Minggu 3.}
\label{fig:peta_jalan_pd_w3}
\end{figure}

\section{Pendahuluan \& Filosofi Fisika Rekayasa: Relaksasi Energi dan Skala Waktu Alami}

Segala fenomena transien dalam rekayasa elektro berpijak pada hukum universal alam: \textbf{alam semesta menolak perubahan energi secara diskontinu (seketika)}. 
\begin{itemize}[leftmargin=*]
    \item Pada rangkaian listrik, daya adalah laju aliran energi: $p(t) = \frac{dE(t)}{dt}$. Agar energi dapat berubah seketika ($\Delta E \neq 0$ dalam interval $\Delta t \to 0$), dibutuhkan daya tak terhingga ($p(t) \to \infty$). Sumber fisis mana pun tidak mampu menyuplai daya tak hingga.
    \item Akibatnya, energi medan magnetik induktor $E_L = \frac{1}{2} L i^2$ dan energi medan elektrostatik kapasitor $E_C = \frac{1}{2} C v^2$ harus bersifat kontinu terhadap waktu.
    \item Demikian pula pada fenomena termal: temperatur suatu transformator sebanding dengan energi termal internal ($E_{\text{th}} = C_{\text{th}} T$). Suhu inti dan minyak trafo tidak dapat melonjak seketika saat terjadi kenaikan arus beban; peningkatan suhu memerlukan waktu untuk menyerap energi kalor dari rugi-rugi tembaga.
\end{itemize}

Oleh karena itu, setiap sistem fisis orde satu memiliki \textbf{skala waktu alami} (\textit{natural time constant}, $\tau$). Konstanta waktu $\tau$ merupakan ukuran inersia sistem dalam merespons eksitasi luar. Semakin besar nilai $\tau$, semakin lambat sistem mencapai kondisi tunak (\textit{steady-state}).

\section{Pemodelan Fisis Transien Rangkaian RL \& Inductive Kickback}

\subsection{Prinsip Dasar Induktor \& Kontinuitas Arus}

Tinjau induktor linear dengan induktansi $L$ Henry. Berdasarkan Hukum Induksi Elektromagnetik Faraday dan Hukum Lenz, perubahan fluks magnetik $\lambda(t) = L i(t)$ membangkitkan gaya gerak listrik (GGL) lawan:
\begin{equation}
    v_L(t) = \frac{d\lambda(t)}{dt} = L \frac{di(t)}{dt}
    \label{eq:inductor_constitutive}
\end{equation}
Persamaan~\eqref{eq:inductor_constitutive} menunjukkan bahwa jika arus berubah seketika dari nilai $I_1$ ke $I_2$ dalam selang waktu $\Delta t \to 0$, maka turunan $\frac{di}{dt} \to \pm \infty$, yang akan menghasilkan tegangan induktif tak terhingga. Karena tegangan tak terhingga menuntut energi tak berhingga dalam waktu sekejap, maka berlaku \textbf{Aksioma Kontinuitas Arus Induktor}:
\begin{equation}
    i_L(0^-) = i_L(0^+)
    \label{eq:inductor_continuity}
\end{equation}

\subsection{Rangkaian RL Pengisian: Respon Langkah (\textit{Step Response})}

Tinjau rangkaian seri pada Gambar~\ref{fig:rl_step_circuit} di mana sebuah induktor $L$ dan resistor $R$ dihubungkan ke sumber tegangan konstan $V_0$ melalui penutupan sakelar pada saat $t = 0$.

\begin{figure}[htbp]
\centering
\begin{circuitikz}[scale=1.0, american]
    \draw (0,0) to[battery2, l=$V_0$] (0,2.5)
          to[closing switch, l={$t=0$}] (2,2.5)
          to[R, l=$R$, v={$v_R(t)$}] (4.5,2.5)
          to[L, l=$L$, v={$v_L(t)$}, i={$i(t)$}] (6.5,2.5)
          -- (6.5,0) -- (0,0);
\end{circuitikz}
\caption{Rangkaian RL seri yang dieksitasi oleh sumber tegangan DC $V_0$ pada $t=0$.}
\label{fig:rl_step_circuit}
\end{figure}

Penerapan Hukum Tegangan Kirchhoff (KVL) pada loop tertutup menghasilkan:
\begin{equation}
    v_R(t) + v_L(t) = V_0 \implies R i(t) + L \frac{di(t)}{dt} = V_0
    \label{eq:rl_kvl}
\end{equation}
Bagi kedua ruas dengan $L$ untuk mendapatkan bentuk standar PDB linier orde 1:
\begin{equation}
    \frac{di(t)}{dt} + \frac{R}{L} i(t) = \frac{V_0}{L}
    \label{eq:rl_standard}
\end{equation}
Selesaikan menggunakan metode faktor pengintegrasi:
\begin{equation}
    P(t) = \frac{R}{L} \implies \mu(t) = \exp\left(\int \frac{R}{L} dt\right) = e^{\frac{R}{L} t}
\end{equation}
Kalikan persamaan standar~\eqref{eq:rl_standard} dengan $\mu(t)$:
\begin{equation}
    e^{\frac{R}{L} t} \frac{di}{dt} + \frac{R}{L} e^{\frac{R}{L} t} i = \frac{d}{dt}\left[ i(t) e^{\frac{R}{L} t} \right] = \frac{V_0}{L} e^{\frac{R}{L} t}
\end{equation}
Integralkan kedua sisi terhadap $t$:
\begin{equation}
    i(t) e^{\frac{R}{L} t} = \int \frac{V_0}{L} e^{\frac{R}{L} t} dt = \frac{V_0}{L} \left( \frac{L}{R} e^{\frac{R}{L} t} \right) + K = \frac{V_0}{R} e^{\frac{R}{L} t} + K
\end{equation}
Kalikan kedua sisi dengan $e^{-\frac{R}{L} t}$:
\begin{equation}
    i(t) = \frac{V_0}{R} + K e^{-\frac{R}{L} t}
    \label{eq:rl_general_sol}
\end{equation}
Terapkan syarat awal: sebelum sakelar ditutup ($t = 0^-$), arus bernilai nol. Sesuai prinsip kontinuitas:
\begin{equation}
    i(0^+) = i(0^-) = 0 \implies \frac{V_0}{R} + K = 0 \implies K = -\frac{V_0}{R}
\end{equation}
Substitusikan nilai $K$ kembali ke persamaan~\eqref{eq:rl_general_sol}:
\begin{equation}
    i(t) = \frac{V_0}{R} \left( 1 - e^{-\frac{t}{\tau}} \right), \quad \text{dengan } \tau = \frac{L}{R} \text{ [detik]}
    \label{eq:rl_step_solution}
\end{equation}
Tegangan pada induktor diperoleh langsung melalui diferensiasi analitik:
\begin{equation}
    v_L(t) = L \frac{di(t)}{dt} = L \left( \frac{V_0}{R} \cdot \frac{1}{\tau} e^{-t/\tau} \right) = L \left( \frac{V_0}{R} \cdot \frac{R}{L} e^{-t/\tau} \right) = V_0 e^{-t/\tau}
    \label{eq:vl_step_solution}
\end{equation}
Perhatikan perilaku batas fisis yang dihasilkan:
\begin{enumerate}[leftmargin=*]
    \item Pada $t = 0^+$ (sesaat setelah sakelar ditutup): $i(0^+) = 0$ dan $v_L(0^+) = V_0$. Induktor berlaku seperti \textbf{rangkaian terbuka} (\textit{open circuit}) sesaat yang menahan seluruh tegangan sumber demi mencegah perubahan arus mendadak.
    \item Pada $t \to \infty$ (kondisi tunak DC): $i(\infty) = \frac{V_0}{R}$ dan $v_L(\infty) = 0$. Induktor berlaku seperti \textbf{hubung singkat} (\textit{short circuit}) sempurna karena arus telah konstan ($\frac{di}{dt} = 0$).
\end{enumerate}

\subsection{Fenomena Bahaya: Lonjakan Tegangan Induktif (\textit{Inductive Kickback})}

Tinjau kondisi ketika rangkaian RL telah mencapai arus tunak $I_0 = \frac{V_0}{R}$. Apabila sakelar tiba-tiba dibuka pada $t = 0$, udara di antara kontak sakelar membentuk celah dengan resistansi yang sangat tinggi ($R_{\text{switch}} \approx 10^7\text{ s.d. } 10^9\,\Omega$), sehingga resistansi total loop melonjak drastis.

Induktor berusaha mempertahankan arus $I_0$ agar tetap mengalir. Sesuai $v_L = L \frac{di}{dt}$, karena arus dipaksa turun ke nol dalam selang waktu pemutusan yang amat sempit ($\Delta t \approx 1\,\mu\text{s}$), laju perubahan arus bernilai negatif yang luar biasa besar ($\frac{di}{dt} \ll 0$). Akibatnya, pada ujung-ujung induktor terbangkit tegangan balik raksasa yang disebut \textbf{inductive kickback}:
\begin{equation}
    v_{\text{kick}} = -L \frac{di}{dt} \gg V_0
    \label{eq:inductive_kickback}
\end{equation}
Tegangan balik ini dapat mencapai ribuan volt pada koil relay 12V/24V, yang memicu:
\begin{enumerate}[leftmargin=*]
    \item \textbf{Busur Api Listrik (\textit{Electric Arcing}):} Mengionisasi udara di antara kontak sakelar mekanik, menyebabkan erosi logam kontak dan kegagalan isolasi.
    \item \textbf{Tembus Tegangan Semikonduktor (\textit{Avalanche Breakdown}):} Menghancurkan transistor pensaklaran (BJT, MOSFET, IGBT) dalam hitungan mikrodetik jika rating $V_{DS}$ terlampaui.
\end{enumerate}

\subsection{Mitigasi Standar Industri: Sirkuit Dioda Freewheeling}

Untuk meniadakan bahaya lonjakan induktif, standar perancangan industri (misalnya pengendali motor DC, katup solenoid, dan relai industri) mewajibkan pemasangan sebuah \textbf{dioda freewheeling} (\textit{flyback diode}) secara paralel-kebalikan (\textit{anti-parallel}) dengan beban induktif, sebagaimana diilustrasikan pada Gambar~\ref{fig:freewheeling_circuit}.

\begin{figure}[htbp]
\centering
\begin{circuitikz}[scale=1.0, american]
    \draw (0,0) to[battery2, l={$V_0$}] (0,3)
          to[opening switch, l={SW ($t=0$)}] (2.5,3)
          to[short, -*] (3.5,3);
    
    % Cabang Paralel Dioda & RL
    \draw (3.5,3) -- (5.5,3)
          to[R, l={$R$}] (5.5,1.5)
          to[L, l={$L$}, i={$i_L(t)$}] (5.5,0)
          to[short, -*] (3.5,0) -- (0,0);
          
    \draw (3.5,3) to[diode, l_={Dioda Flyback $D$}] (3.5,0);
\end{circuitikz}
\caption{Sirkuit proteksi induktif menggunakan dioda \textit{freewheeling} anti-paralel.}
\label{fig:freewheeling_circuit}
\end{figure}

\textbf{Mekanisme Operasi Transien:}
\begin{itemize}[leftmargin=*]
    \item \textbf{Saat Sakelar Tertutup ($t < 0$):} Katoda dioda terhubung ke potensial positif $V_0$ dan anoda ke ground. Dioda berada pada kondisi panjar mundur (\textit{reverse-biased}) dan bertindak sebagai sirkit terbuka. Arus mengalir normal melewati loop $R-L$.
    \item \textbf{Saat Sakelar Dibuka ($t = 0$):} Arus induktor tidak dapat terputus seketika. GGL lawan membalik polaritas tegangan induktor (ujung bawah menjadi positif terhadap ujung atas). Hal ini langsung memicu dioda berada pada panjar maju (\textit{forward-biased}).
    \item Arus induktor dialihkan bersirkulasi secara aman di dalam loop tertutup dioda-resistor-induktor. Energi medan magnetik $\frac{1}{2} L I_0^2$ diserap dan didisipasikan secara eksponensial menjadi panas oleh resistansi kawat koil dan resistansi internal dioda:
    \begin{equation}
        i_L(t) = I_0 e^{-\frac{R + R_D}{L} t}
        \label{eq:freewheeling_decay}
    \end{equation}
    Tegangan maksimum yang dialami sakelar terkunci pada nilai yang sangat aman: $V_{\text{max}} = V_0 + V_{\text{forward-diode}} \approx V_0 + 0{,}7\text{ Volt}$.
\end{itemize}

\section{Pemodelan Fisis Transien Rangkaian RC \& Analisis Multi-Interval}

\subsection{Respon Alami \& Respon Langkah Kapasitor}

Pada Modul 2, telah dibuktikan bahwa tegangan kapasitor pada rangkaian RC memenuhi PDB linier orde 1:
\begin{equation}
    \frac{dv_C(t)}{dt} + \frac{1}{RC} v_C(t) = \frac{v_s(t)}{RC}
    \label{eq:rc_standard_diff}
\end{equation}
Solusi umum untuk eksitasi tegangan konstan $V_0$ dengan syarat awal $v_C(0) = V_{\text{init}}$ adalah:
\begin{equation}
    v_C(t) = V_0 + (V_{\text{init}} - V_0) e^{-t/\tau}, \quad \tau = RC
    \label{eq:rc_complete_response}
\end{equation}
Formula~\eqref{eq:rc_complete_response} dikenal sebagai \textbf{Formula Respon Lengkap Tiga Parameter}:
\begin{equation}
    v(t) = v(\infty) + \left[ v(0^+) - v(\infty) \right] e^{-t/\tau}
    \label{eq:three_parameter_formula}
\end{equation}
di mana:
\begin{itemize}[leftmargin=*]
    \item $v(0^+)$ adalah nilai awal saat transien dimulai (ditentukan oleh kontinuitas energi).
    \item $v(\infty)$ adalah nilai akhir kondisi tunak setelah seluruh transien meluruh.
    \item $\tau$ adalah konstanta waktu ekivalen Thévenin yang dilihat oleh elemen penyimpan energi: $\tau = R_{\text{th}} C$ atau $\tau = L / R_{\text{th}}$.
\end{itemize}

\subsection{Analisis Transien Multi-Interval (Switching Dinamis)}

Dalam aplikasi sistem proteksi dan pembangkit pulsa radar/telekomunikasi, sakelar beroperasi beberapa kali pada waktu-waktu yang berbeda ($t_1, t_2, \dots$). Kunci penyelesaian PDB multi-interval adalah \textbf{kontinuitas variabel keadaan pada batas interval}:
\begin{equation}
    v_C(t_1^+) = v_C(t_1^-) \quad \text{dan} \quad i_L(t_1^+) = i_L(t_1^-)
\end{equation}
Setiap interval waktu dimodelkan sebagai IVP baru dengan titik acuan waktu digeser sebesar $(t - t_k)$:
\begin{equation}
    v(t) = v(\infty)_k + \left[ v(t_k^+) - v(\infty)_k \right] e^{-\frac{t - t_k}{\tau_k}}, \quad \text{untuk } t_k \le t < t_{k+1}
\end{equation}

\section{Pemodelan Termal Transformator Berbasis Hukum Pendinginan Newton}

\subsection{Filosofi Perpindahan Kalor pada Transformator Daya}

Transformator daya pada gardu induk PLN menyalurkan daya puluhan hingga ratusan Megavolt-Ampere (MVA). Efisiensi transformator sangat tinggi ($\approx 98{,}5\%$ s.d. $99{,}5\%$), namun sisa $0{,}5\%$ s.d. $1{,}5\%$ disipasi daya diubah menjadi panas di dalam tangki transformator. Pada trafo 60 MVA, disipasi $1\%$ menghasilkan daya rugi-rugi kalor sebesar:
\begin{equation}
    P_{\text{loss}} = 0{,}01 \times 60 \times 10^6\text{ W} = 600\text{ kW!}
\end{equation}
Daya kalor 600 kW setara dengan pemanas air listrik skala raksasa yang terus-menerus memanaskan minyak dielektrik dan belitan tembaga. Jika panas ini tidak dipindahkan ke udara lingkungan, suhu internal akan melonjak melampaui batas ketahanan termal bahan isolasi kertas selulosa ($105^\circ\text{C}$ s.d. $140^\circ\text{C}$), menyebabkan pembentukan gas terlarut (\textit{Dissolved Gas Analysis} / DGA) dan ledakan tangki trafo.

\subsection{Penurunan Persamaan Diferensial Termal dari Prinsip Pertama}

Hukum Kekekalan Energi Termal (Hukum Pertama Termodinamika) untuk volume kontrol transformator menyatakan:
\begin{equation}
\begin{split}
    \left[ \text{Laju Kalor Dihasilkan} \right] &- \left[ \text{Laju Kalor Dibuang ke Lingkungan} \right] \\
    &= \left[ \text{Laju Akumulasi Energi Kalor Internal} \right]
\end{split}
\label{eq:thermal_balance}
\end{equation}
Definisikan variabel-variabel termal sistem:
\begin{itemize}[leftmargin=*]
    \item $T(t)$ : Temperatur rata-rata minyak/inti transformator pada waktu $t$ [$^\circ$C].
    \item $T_{\text{amb}}$ : Temperatur udara ambien lingkungan sekitar tangki trafo [$^\circ$C].
    \item $P_{\text{loss}}$ : Rugi-rugi daya termal total ($P_{\text{loss}} = P_{\text{inti}} + I^2 R_{\text{belitan}}$) [Watt].
    \item $C_{\text{th}}$ : Kapasitas panas efektif massa total transformator (besi inti + tembaga belitan + minyak isolasi), dalam satuan $[\text{Joule}/^\circ\text{C}]$.
    \item $R_{\text{th}}$ : Resistansi termal konveksi dan radiasi sirip radiator pendingin, dalam satuan $[^\circ\text{C}/\text{Watt}]$.
    \item $h$ : Koefisien konveksi permukaan, $A$ : Luas penampang pendingin radiator tangki, di mana $R_{\text{th}} = \frac{1}{h A}$.
\end{itemize}

Berdasarkan \textbf{Hukum Pendinginan Newton}, laju pembuangan panas ke lingkungan berbanding lurus dengan selisih temperatur antara permukaan pendingin dan ambien:
\begin{equation}
    q_{\text{buang}}(t) = h A \left( T(t) - T_{\text{amb}} \right) = \frac{T(t) - T_{\text{amb}}}{R_{\text{th}}}
    \label{eq:newton_cooling}
\end{equation}
Laju akumulasi energi panas di dalam massa transformator diberikan oleh:
\begin{equation}
    q_{\text{akumulasi}}(t) = C_{\text{th}} \frac{dT(t)}{dt}
    \label{eq:thermal_accumulation}
\end{equation}
Substitusikan persamaan~\eqref{eq:newton_cooling} dan~\eqref{eq:thermal_accumulation} ke neraca energi~\eqref{eq:thermal_balance}:
\begin{equation}
    P_{\text{loss}} - \frac{T(t) - T_{\text{amb}}}{R_{\text{th}}} = C_{\text{th}} \frac{dT(t)}{dt}
    \label{eq:thermal_ode_raw}
\end{equation}
Bagi kedua ruas dengan $C_{\text{th}}$ dan atur ke dalam bentuk kanonik PDB linier orde 1:
\begin{equation}
    \frac{dT(t)}{dt} + \frac{1}{R_{\text{th}} C_{\text{th}}} T(t) = \frac{T_{\text{amb}}}{R_{\text{th}} C_{\text{th}}} + \frac{P_{\text{loss}}}{C_{\text{th}}}
    \label{eq:thermal_ode_standard}
\end{equation}

\subsection{Analogi Rangkaian Listrik-Termal (Lumped Parameter Model)}

Persamaan diferensial~\eqref{eq:thermal_ode_standard} identik secara matematis dengan persamaan KCL rangkaian RC paralel yang disuplai oleh sumber arus independen, seperti dipaparkan pada Tabel~\ref{tab:analogi_elektro_termal}.

\begin{table}[htbp]
\centering
\caption{Analogi Fisis antara Variabel Sistem Elektrik dan Sistem Termal.}
\label{tab:analogi_elektro_termal}
\resizebox{\linewidth}{!}{%
\begin{tabular}{llll}
\toprule
\textbf{Besaran Sistem} & \textbf{Domain Rangkaian Listrik} & \textbf{Domain Termal Trafo} & \textbf{Hubungan Konstitutif Fisis} \\ \midrule
Potensial Penggerak & Beda Tegangan $v(t)$ [V] & Beda Temperatur $\theta(t) = T - T_{\text{amb}}$ [$^\circ$C] & Hukum Ohm / Hukum Fourier \\
Laju Aliran Fluks & Arus Listrik $i(t)$ [A] & Daya Rugi Kalor $P_{\text{loss}}(t)$ [W] & Laju perpindahan muatan vs kalor \\
Parameter Disipatif & Resistansi $R$ [$\Omega$] & Resistansi Termal $R_{\text{th}}$ [$^\circ$C/W] & Disipasi energi Joule vs konveksi \\
Parameter Penyimpan & Kapasitansi $C$ [Farad] & Kapasitas Panas $C_{\text{th}}$ [J/$^\circ$C] & Penyimpanan muatan vs energi dalam \\
Konstanta Waktu & $\tau = R C$ [detik] & $\tau_{\text{th}} = R_{\text{th}} C_{\text{th}}$ [menit/jam] & Laju relaksasi eksponensial \\ \bottomrule
\end{tabular}%
}
\end{table}

\subsection{Derivasi Analitis Solusi Respon Temperatur Transformator}

Definisikan konstanta waktu termal:
\begin{equation}
    \tau_{\text{th}} = R_{\text{th}} C_{\text{th}}
    \label{eq:thermal_time_constant}
\end{equation}
Persamaan diferensial~\eqref{eq:thermal_ode_standard} dapat ditulis ulang dalam suku kenaikan suhu di atas ambien, $\theta(t) = T(t) - T_{\text{amb}}$:
\begin{equation}
    \frac{d\theta(t)}{dt} + \frac{1}{\tau_{\text{th}}} \theta(t) = \frac{P_{\text{loss}}}{C_{\text{th}}} = \frac{R_{\text{th}} P_{\text{loss}}}{\tau_{\text{th}}}
\end{equation}
Pada kondisi tunak ($t \to \infty$), laju perubahan suhu mendekati nol ($\frac{d\theta}{dt} = 0$), sehingga kenaikan suhu tunak maksimum ($\Delta T_{\text{ss}}$) adalah:
\begin{equation}
    \theta(\infty) = \Delta T_{\text{ss}} = P_{\text{loss}} \cdot R_{\text{th}} \quad [^\circ\text{C}]
    \label{eq:steady_state_temp_rise}
\end{equation}
Dengan faktor integrasi $\mu(t) = e^{t/\tau_{\text{th}}}$ dan kondisi awal saat $t = 0$ yaitu $T(0) = T_0$, solusi analitis lengkap untuk temperatur transformator setiap saat adalah:
\begin{equation}
    T(t) = T_{\text{amb}} + \Delta T_{\text{ss}} \left( 1 - e^{-\frac{t}{\tau_{\text{th}}}} \right) + (T_0 - T_{\text{amb}}) e^{-\frac{t}{\tau_{\text{th}}}}
    \label{eq:thermal_complete_solution}
\end{equation}

\subsection{Standar Industri IEEE Std C57.91: Penuaan Isolasi \& Pembebanan Darurat}

Standar internasional \textbf{IEEE Std C57.91} (\textit{Guide for Loading Mineral-Oil-Immersed Transformers}) mendasarkan masa pakai transformator pada degradasi termal kekuatan mekanik kertas insulasi selulosa. 
Laju penuaan relatif insulasi (\textit{Relative Aging Rate}, $V$) dihitung menggunakan persamaan reaksi kimiawi Arrhenius:
\begin{equation}
    V = \exp\left( \frac{15000}{383} - \frac{15000}{\Theta_{\text{H}} + 273} \right) \approx 2^{\frac{\Theta_{\text{H}} - 110}{6}}
    \label{eq:arrhenius_aging}
\end{equation}
di mana $\Theta_{\text{H}}$ adalah temperatur titik terpanas belitan (\textit{winding hot-spot temperature}) dalam $^\circ$C.
\begin{itemize}[leftmargin=*]
    \item Pada temperatur acuan nominal $\Theta_{\text{H}} = 110^\circ\text{C}$, nilai penuaan relatif adalah $V = 1{,}0$ (umur operasi normal trafo didesain selama 180.000 jam atau $\approx 20{,}55$ tahun).
    \item Setiap kenaikan temperatur titik terpanas sebesar $6^\circ\text{C}$ di atas $110^\circ\text{C}$ melipatgandakan laju penuaan insulasi menjadi dua kali lipat ($V = 2$ pada $116^\circ\text{C}$, $V = 4$ pada $122^\circ\text{C}$, dan $V = 8$ pada $128^\circ\text{C}$).
    \item Sebaliknya, jika sistem pendingin trafo bekerja optimal sehingga suhu berada di bawah $110^\circ\text{C}$, laju penuaan melambat secara signifikan ($V < 1{,}0$).
\end{itemize}

Oleh karena itu, penyelesaian PDB termal secara waktu nyata (\textit{real-time dynamic thermal rating}) sangat krusial bagi operator gardu induk PLN untuk menentukan berapa lama trafo dapat memikul beban lebih darurat (\textit{emergency overload}) tanpa merusak insulasi internalnya.

\section{Contoh Soal Terhitung Numerik Langkah demi Langkah}

\begin{example}[\textbf{Transien Pengisian RL \& Penentuan Dioda Flyback Sakelar DC}]
Sebuah koil kontaktor magnetik pada panel distribusi industri memiliki induktansi $L = 2{,}5\text{ H}$ dan resistansi internal kawat koil $R = 50\,\Omega$. Koil ini disuplai oleh sumber DC $V_0 = 24\text{ V}$ melalui sakelar transistor MOSFET pada $t = 0$.
\begin{enumerate}[label=(\alph*)]
    \item Tentukan konstanta waktu rangkaian $\tau$ dan rumus matematis arus koil $i(t)$.
    \item Hitung arus dan energi medan magnetik yang tersimpan pada kondisi tunak ($t \to \infty$).
    \item Hitung waktu yang dibutuhkan arus koil untuk mencapai $90\%$ dari nilai maksimumnya.
    \item Jika saat kondisi tunak sakelar tiba-tiba dibuka tanpa dioda freewheeling dan busur api memutus arus dalam waktu $\Delta t = 2\,\mu\text{s}$, estimasikan nilai puncak tegangan balik induktif (\textit{inductive kickback}).
    \item Jika dipasang dioda silikon freewheeling ($V_F = 0{,}7\text{ V}$), tentukan tegangan maksimum yang dialami transistor.
\end{enumerate}

\paragraph{Penyelesaian Langkah demi Langkah:}
\begin{enumerate}[label=(\alph*), leftmargin=*]
    \item \textbf{Konstanta waktu listrik rangkaian RL:}
    \begin{equation*}
        \tau = \frac{L}{R} = \frac{2{,}5\text{ H}}{50\,\Omega} = 0{,}05\text{ detik} = 50\text{ ms}
    \end{equation*}
    Arus tunak akhir:
    \begin{equation*}
        I_{\text{maks}} = \frac{V_0}{R} = \frac{24\text{ V}}{50\,\Omega} = 0{,}48\text{ A}
    \end{equation*}
    Substitusikan ke solusi khusus respon langkah:
    \begin{equation*}
        i(t) = 0{,}48 \left( 1 - e^{-\frac{t}{0{,}05}} \right) = 0{,}48 \left( 1 - e^{-20 t} \right)\text{ [Ampere]}
    \end{equation*}

    \item \textbf{Energi medan magnetik pada kondisi tunak:}
    \begin{equation*}
        E_L = \frac{1}{2} L I_{\text{maks}}^2 = \frac{1}{2} \times 2{,}5\text{ H} \times (0{,}48\text{ A})^2 = 1{,}25 \times 0{,}2304 = 0{,}288\text{ Joule}
    \end{equation*}

    \item \textbf{Waktu mencapai 90\% nilai tunak ($i(t_1) = 0{,}90 \times 0{,}48 = 0{,}432\text{ A}$):}
    \begin{align*}
        0{,}48 \left( 1 - e^{-t_1/\tau} \right) = 0{,}432 &\implies 1 - e^{-t_1/\tau} = 0{,}90 \\
        &\implies e^{-t_1/\tau} = 0{,}10 \\
        &\implies -\frac{t_1}{\tau} = \ln(0{,}10) = -2{,}3026 \\
        &\implies t_1 = 2{,}3026 \times 0{,}05\text{ s} \approx 0{,}1151\text{ s} = 115{,}1\text{ ms}
    \end{align*}

    \item \textbf{Estimasi lonjakan tegangan balik tanpa proteksi:}
    Laju penurunan arus rata-rata saat kontak terbuka dalam $\Delta t = 2\,\mu\text{s} = 2 \times 10^{-6}\text{ s}$:
    \begin{equation*}
        \left|\frac{\Delta i}{\Delta t}\right| \approx \frac{0{,}48\text{ A} - 0\text{ A}}{2 \times 10^{-6}\text{ s}} = 240.000\text{ A/s}
    \end{equation*}
    Besar tegangan balik induktif yang dibangkitkan:
    \begin{equation*}
        |v_L| = L \left| \frac{\Delta i}{\Delta t} \right| = 2{,}5\text{ H} \times 240.000\text{ A/s} = 600.000\text{ Volt!}
    \end{equation*}
    \textit{Catatan Rekayasa:} Tegangan teoretis $600\text{ kV}$ ini dalam kenyataannya akan langsung memicu tembus listrik (\textit{spark breakdown}) pada celah sakelar atau menghancurkan transistor semikonduktor dengan rating breakdown tipikal $50\text{ s.d. } 100\text{ V}$.

    \item \textbf{Tegangan maksimum pada sakelar dengan dioda freewheeling:}
    Ketika sakelar dibuka, anoda dioda berada pada ground ($0\text{ V}$) dan katoda tersambung ke titik kolektor/drain transistor. Dioda membatasi tegangan jatuh maju sebesar $V_F = 0{,}7\text{ V}$:
    \begin{equation*}
        V_{\text{transistor, maks}} = V_0 + V_F = 24\text{ V} + 0{,}7\text{ V} = 24{,}7\text{ Volt}
    \end{equation*}
    Tegangan ini berada jauh di bawah rating aman transistor ($100\text{ V}$), sehingga sirkuit terlindungi secara sempurna. \qed
\end{enumerate}
\end{example}

\begin{example}[\textbf{Transien Rangkaian RC Switching Multi-Interval}]
Tinjau sirkuit RC dengan parameter $C = 100\,\mu\text{F}$, $R_1 = 20\text{ k}\Omega$, dan $R_2 = 5\text{ k}\Omega$. Sakelar berada pada posisi A (pengisian dari sumber $V_0 = 100\text{ V}$) untuk selang waktu $0 \le t < 4\text{ detik}$, kemudian dipindahkan ke posisi B (pengosongan melalui $R_2$) pada saat $t = 4\text{ detik}$. Kapasitor awalnya tidak bermuatan ($v_C(0) = 0$).

\begin{center}
\begin{circuitikz}[scale=0.95, american]
    \draw (0,0) to[battery2, l={$V_0=100\text{ V}$}] (0,2.5)
          -- (1.5,2.5) node[above] {A};
    \draw (1.5,1.0) node[below] {B} -- (2.5,1.0) to[R, l={$R_2=5\text{ k}\Omega$}] (2.5,0) -- (0,0);
    \draw (2.2,1.75) node[left] {SW} to[short] (3.0,1.75)
          to[R, l={$R_1=20\text{ k}\Omega$}] (5.2,1.75)
          to[C, l={$C=100\,\mu\text{F}$}, v={$v_C(t)$}] (5.2,0)
          -- (2.5,0);
    \draw[dashed, ->] (2.2,1.75) -- (1.6,2.4);
    \draw[dashed, ->] (2.2,1.75) -- (1.6,1.1);
\end{circuitikz}
\end{center}

Tentukan:
\begin{enumerate}[label=(\alph*)]
    \item Tegangan kapasitor $v_C(t)$ pada interval pengisian $0 \le t < 4\text{ detik}$.
    \item Nilai tegangan kapasitor tepat sebelum sakelar dipindahkan ($t = 4^-\text{ s}$).
    \item Tegangan kapasitor $v_C(t)$ pada interval pengosongan $t \ge 4\text{ detik}$.
    \item Waktu total $t$ (dihitung sejak $t=0$) saat tegangan kapasitor turun menjadi $10\text{ V}$.
\end{enumerate}

\paragraph{Penyelesaian Langkah demi Langkah:}
\begin{enumerate}[label=(\alph*), leftmargin=*]
    \item \textbf{Interval 1: $0 \le t < 4\text{ s}$ (Posisi A):}
    Resistansi pengisian adalah $R_1 = 20\text{ k}\Omega = 20 \times 10^3\,\Omega$. Konstanta waktu:
    \begin{equation*}
        \tau_1 = R_1 C = (20 \times 10^3\,\Omega) \times (100 \times 10^{-6}\text{ F}) = 2{,}0\text{ detik}
    \end{equation*}
    Karena $v_C(0) = 0$ dan target tunak $v(\infty) = 100\text{ V}$, solusi pengisian adalah:
    \begin{equation*}
        v_C(t) = 100 \left( 1 - e^{-t/2} \right)\text{ [Volt]}, \quad 0 \le t < 4\text{ s}
    \end{equation*}

    \item \textbf{Tegangan pada $t = 4^-\text{ s}$:}
    \begin{equation*}
        v_C(4^-) = 100 \left( 1 - e^{-4/2} \right) = 100 \left( 1 - e^{-2} \right) = 100 \left( 1 - 0{,}1353 \right) = 86{,}47\text{ Volt}
    \end{equation*}

    \item \textbf{Interval 2: $t \ge 4\text{ s}$ (Posisi B):}
    Prinsip kontinuitas tegangan menjamin:
    \begin{equation*}
        v_C(4^+) = v_C(4^-) = 86{,}47\text{ Volt}
    \end{equation*}
    Pada posisi B, sumber tegangan terputus ($V_{\text{sumber}} = 0$) dan kapasitor terhubung seri dengan loop resistor $R_1$ dan $R_2$. Resistansi total loop pengosongan adalah:
    \begin{equation*}
        R_{\text{total}} = R_1 + R_2 = 20\text{ k}\Omega + 5\text{ k}\Omega = 25\text{ k}\Omega
    \end{equation*}
    Konstanta waktu pengosongan baru:
    \begin{equation*}
        \tau_2 = R_{\text{total}} C = (25 \times 10^3\,\Omega) \times (100 \times 10^{-6}\text{ F}) = 2{,}5\text{ detik}
    \end{equation*}
    Solusi umum pengosongan untuk waktu acuan yang digeser $(t - 4)$:
    \begin{equation*}
        v_C(t) = v_C(4^+) e^{-\frac{t - 4}{\tau_2}} = 86{,}47 \, e^{-\frac{t - 4}{2{,}5}}\text{ [Volt]}, \quad t \ge 4\text{ s}
    \end{equation*}

    \item \textbf{Waktu saat $v_C(t) = 10\text{ V}$:}
    Peristiwa ini terjadi pada interval pengosongan ($t \ge 4$):
    \begin{align*}
        86{,}47 \, e^{-\frac{t - 4}{2{,}5}} = 10 &\implies e^{-\frac{t - 4}{2{,}5}} = \frac{10}{86{,}47} = 0{,}11565 \\
        &\implies -\frac{t - 4}{2{,}5} = \ln(0{,}11565) = -2{,}1572 \\
        &\implies t - 4 = 2{,}5 \times 2{,}1572 = 5{,}393\text{ detik} \\
        &\implies t = 4 + 5{,}393 = 9{,}393\text{ detik}
    \end{align*}
    Jadi, tegangan kapasitor meluruh menjadi $10\text{ V}$ pada $t \approx 9{,}39\text{ detik}$ sejak awal pengisian. \qed
\end{enumerate}
\end{example}

\begin{example}[\textbf{Manajemen Termal Transformator Daya Gardu Induk 60 MVA}]
Sebuah transformator daya 60 MVA 150/20 kV pada gardu induk PLN mengalami lonjakan beban industri saat jam beban puncak (\textit{peak load}). Parameter termal transformator adalah:
\begin{itemize}[leftmargin=*]
    \item Kapasitas panas total transformator: $C_{\text{th}} = 3{,}6 \times 10^6\text{ J}/^\circ\text{C}$.
    \item Resistansi termal pendingin radiator: $R_{\text{th}} = 1{,}25 \times 10^{-4}\;^\circ\text{C}/\text{W}$.
    \item Temperatur udara lingkungan (ambien): $T_{\text{amb}} = 30^\circ\text{C}$.
    \item Temperatur awal minyak sebelum lonjakan beban ($t=0$): $T_0 = 55^\circ\text{C}$.
    \item Total rugi-rugi daya saat beban puncak (rugi inti + tembaga): $P_{\text{loss}} = 480\text{ kW} = 480.000\text{ W}$.
\end{itemize}
Standar operasi gardu induk menetapkan alarm peringatan suhu minyak atas pada $T_{\text{alarm}} = 85^\circ\text{C}$ dan batas trip pemutus tenaga (CB) pada $T_{\text{trip}} = 95^\circ\text{C}$.
\begin{enumerate}[label=(\alph*)]
    \item Hitung konstanta waktu termal transformator $\tau_{\text{th}}$ dalam satuan menit dan jam.
    \item Tentukan temperatur tunak akhir minyak transformator ($T_{\text{ss}}$) jika pembebanan puncak dibiarkan berlangsung terus-menerus. Apakah transformator akan trip?
    \item Turunkan persamaan fungsi temperatur minyak transformator terhadap waktu $T(t)$.
    \item Hitung berapa lama (dalam jam atau menit) operator gardu induk memiliki waktu untuk memindahkan sebagian beban sebelum alarm peringatan ($85^\circ\text{C}$) berbunyi.
\end{enumerate}

\paragraph{Penyelesaian Langkah demi Langkah:}
\begin{enumerate}[label=(\alph*), leftmargin=*]
    \item \textbf{Konstanta waktu termal transformator:}
    \begin{equation*}
        \tau_{\text{th}} = R_{\text{th}} C_{\text{th}} = (1{,}25 \times 10^{-4}\;^\circ\text{C}/\text{W}) \times (3{,}6 \times 10^6\text{ J}/^\circ\text{C}) = 450\text{ detik}
    \end{equation*}
    Konversi ke menit dan jam:
    \begin{equation*}
        \tau_{\text{th}} = \frac{450}{60} = 7{,}5\text{ menit} = 0{,}125\text{ jam}
    \end{equation*}
    \textit{Catatan Pemodelan:} Nilai $\tau_{\text{th}}$ ini mewakili konstanta waktu termal lokal minyak radiator; transformator skala besar secara utuh memiliki konstanta waktu 2 hingga 4 jam, namun sirip pendingin merespons lebih cepat.

    \item \textbf{Temperatur tunak akhir minyak transformator ($T_{\text{ss}}$):}
    Kenaikan temperatur tunak akibat disipasi panas 480 kW:
    \begin{equation*}
        \Delta T_{\text{ss}} = P_{\text{loss}} \times R_{\text{th}} = 480.000\text{ W} \times (1{,}25 \times 10^{-4}\;^\circ\text{C}/\text{W}) = 60^\circ\text{C}
    \end{equation*}
    Temperatur mutlak tunak:
    \begin{equation*}
        T_{\text{ss}} = T_{\text{amb}} + \Delta T_{\text{ss}} = 30^\circ\text{C} + 60^\circ\text{C} = 90^\circ\text{C}
    \end{equation*}
    \textbf{Evaluasi Keamanan Sistem:}
    Karena $T_{\text{ss}} = 90^\circ\text{C} < T_{\text{trip}} = 95^\circ\text{C}$, pemutus tenaga (CB) tidak akan trip jika kondisi tunak tercapai. Namun, karena $T_{\text{ss}} = 90^\circ\text{C} > T_{\text{alarm}} = 85^\circ\text{C}$, sistem monitoring gardu induk \textbf{pasti akan membunyikan alarm} peringatan beban lebih.

    \item \textbf{Persamaan temperatur terhadap waktu $T(t)$:}
    Gunakan formula umum respon termal orde satu~\eqref{eq:thermal_complete_solution} dengan $T(0) = 55^\circ\text{C}$ dan $T_{\text{ss}} = 90^\circ\text{C}$:
    \begin{align*}
        T(t) &= T_{\text{ss}} + \left[ T(0) - T_{\text{ss}} \right] e^{-t/\tau_{\text{th}}} \\
             &= 90 + (55 - 90) e^{-t/7{,}5} \\
             &= 90 - 35 \, e^{-t/7{,}5}\text{ [}^\circ\text{C]}, \quad \text{dengan } t \text{ dalam menit}
    \end{align*}

    \item \textbf{Waktu hingga alarm $85^\circ\text{C}$ berbunyi ($T(t_{\text{alarm}}) = 85^\circ\text{C}$):}
    \begin{align*}
        90 - 35 \, e^{-t_{\text{alarm}}/7{,}5} = 85 &\implies 35 \, e^{-t_{\text{alarm}}/7{,}5} = 5 \\
        &\implies e^{-t_{\text{alarm}}/7{,}5} = \frac{5}{35} = \frac{1}{7} \approx 0{,}14286 \\
        &\implies -\frac{t_{\text{alarm}}}{7{,}5} = \ln\left(\frac{1}{7}\right) = -\ln(7) = -1{,}9459 \\
        &\implies t_{\text{alarm}} = 7{,}5 \times 1,9459 \approx 14{,}59\text{ menit}
    \end{align*}
    Operator gardu induk memiliki jendela waktu operasional darurat sekitar \textbf{14 menit 35 detik} sejak terjadinya lonjakan beban untuk menyalakan kipas pendingin tambahan (forced-cooling ONAF) atau memindahkan sebagian aliran daya ke penyulang cadangan sebelum alarm berbunyi. \qed
\end{enumerate}
\end{example}

\section{Praktikum Komputasi Numerik Terbuka Berbasis Julia}

Sebagai implementasi Pilar 3 Pedagogis, simulasi komputasi dilakukan menggunakan bahasa ilmiah modern berkecepatan tinggi \textbf{Julia 1.12+} dengan ekosistem \texttt{DifferentialEquations.jl} dan \texttt{Plots.jl}.

\subsection{Skrip Komputasi: Simulasi Transien Elektrik \& Termal Trafo}

Listing~\ref{lst:julia_transient_sim} menyajikan program Julia lengkap untuk memodelkan:
\begin{enumerate}[leftmargin=*]
    \item Transien pengisian dan pembukaan induktor RL dengan mitigasi dioda freewheeling.
    \item Respon kenaikan suhu transformator daya gardu induk dari kondisi beban normal ke beban lebih puncak (overload 130\%).
\end{enumerate}

\begin{lstlisting}[style=juliastyle, caption={Skrip Julia untuk Pemodelan Transien Elektromagnetik dan Dinamika Termal Trafo.}, label={lst:julia_transient_sim}]
# =================================================================
# PRAKTIKUM PERSAMAAN DIFERENSIAL - MINGGU 3
# Pemodelan Transien Elektrik RL & Dinamika Termal Trafo Daya
# Jurusan Teknik Elektro, Universitas Bengkulu
# =================================================================

using DifferentialEquations
using Plots

# -----------------------------------------------------------------
# BAGIAN 1: Pemodelan Transien Rangkaian RL DC
# PDB: di/dt = (V0 - R*i) / L
# -----------------------------------------------------------------
function rl_circuit!(di, i, p, t)
    V0, R, L = p
    di[1] = (V0 - R * i[1]) / L
end

# Parameter Fisik Rangkaian RL
V0 = 24.0          # Tegangan sumber DC [Volt]
R  = 50.0          # Resistansi koil [Ohm]
L  = 2.5           # Induktansi koil [Henry]
tau_el = L / R     # Konstanta waktu listrik [detik] (0.05 s)

p_rl = [V0, R, L]
i0 = [0.0]         # Syarat awal: i(0) = 0 A
tspan_rl = (0.0, 0.25) # Simulasi selama 5 konstanta waktu

prob_rl = ODEProblem(rl_circuit!, i0, tspan_rl, p_rl)
sol_rl = solve(prob_rl, Tsit5(), reltol=1e-8, abstol=1e-8)

# Perhitungan Solusi Eksak Analitik untuk Validasi
t_eval = range(0.0, 0.25, length=200)
i_exact = (V0 / R) .* (1.0 .- exp.(-t_eval ./ tau_el))
i_num = [sol_rl(t)[1] for t in t_eval]
error_max_rl = maximum(abs.(i_num .- i_exact))

println("=== VALIDASI SIMULASI TRANSIEN RL ===")
println("Konstanta waktu listrik (tau) : ", tau_el, " detik")
println("Arus kondisi tunak (I_max)   : ", V0 / R, " Ampere")
println("Galat mutlak maksimum analitik vs numerik : ", error_max_rl, " A")

# -----------------------------------------------------------------
# BAGIAN 2: Pemodelan Termal Transformator Daya Gardu Induk
# PDB: dT/dt = (Tamb - T) / tau_th + (Ploss / Cth)
# -----------------------------------------------------------------
function transformer_thermal!(dT, T, p, t)
    Tamb, Ploss, Cth, Rth = p
    tau_th = Rth * Cth
    dT[1] = -(T[1] - Tamb) / tau_th + (Ploss / Cth)
end

# Parameter Fisis Transformator Daya 60 MVA
Tamb  = 30.0               # Suhu ambien [C]
Ploss = 480_000.0          # Daya rugi-rugi termal total [W]
Cth   = 3.6e6              # Kapasitas panas trafo [J/C]
Rth   = 1.25e-4            # Resistansi termal radiator [C/W]
tau_th_min = (Rth * Cth) / 60.0 # Konstanta waktu [menit]

p_th = [Tamb, Ploss, Cth, Rth]
T0 = [55.0]                # Suhu awal minyak trafo [C]
tspan_th = (0.0, 45.0)     # Simulasi selama 45 menit

prob_th = ODEProblem(transformer_thermal!, T0, tspan_th, p_th)
sol_th = solve(prob_th, Tsit5(), reltol=1e-8, abstol=1e-8)

println("\n=== VALIDASI SIMULASI TERMAL TRANSFORMATOR ===")
println("Konstanta waktu termal (tau_th): ", tau_th_min, " menit")
println("Temperatur tunak akhir (T_ss)  : ", Tamb + Ploss * Rth, " C")

# -----------------------------------------------------------------
# BAGIAN 3: Visualisasi Grafis Respon Dinamis Sistem
# -----------------------------------------------------------------
p1 = plot(t_eval .* 1000, i_exact, label="Analitik Eksak", lw=2.5, color=:blue)
plot!(p1, t_eval .* 1000, i_num, label="Numerik Tsit5()", lw=1.5, ls=:dash, color=:red)
xlabel!(p1, "Waktu t [milidetik]")
ylabel!(p1, "Arus Koil i(t) [Ampere]")
title!(p1, "Transien Pengisian Koil RL Seri (24 V DC)")

t_th_eval = range(0.0, 45.0, length=200)
T_num = [sol_th(t)[1] for t in t_th_eval]

p2 = plot(t_th_eval, T_num, label="Suhu Minyak T(t)", lw=2.5, color=:darkorange)
hline!(p2, [85.0], label="Ambang Alarm (85 C)", lw=1.8, ls=:dash, color=:red)
hline!(p2, [95.0], label="Ambang Trip CB (95 C)", lw=1.8, ls=:dot, color=:black)
xlabel!(p2, "Waktu Operasi [Menit]")
ylabel!(p2, "Temperatur Minyak [C]")
title!(p2, "Dinamika Kenaikan Suhu Trafo Daya 60 MVA")

plot(p1, p2, layout=(2,1), size=(800, 700))
savefig("transien_elektro_termal.png")
println("Simulasi selesai. Grafik disimpan ke transien_elektro_termal.png")
\end{lstlisting}

\section{Evaluasi \& Bank Soal Terstruktur Taksonomi Bloom (C1 s.d. C6)}

Untuk menguji pemahaman mahasiswa secara komprehensif, berikut disajikan evaluasi bertingkat mengacu pada standar Taksonomi Bloom:

\begin{enumerate}[leftmargin=*]
    \item \textbf{[C1 - Mengingat]} Tuliskan hubungan konstitutif tegangan-arus pada induktor dan kapasitor serta jelaskan mengapa arus induktor $i_L(t)$ dan tegangan kapasitor $v_C(t)$ tidak boleh bernilai diskontinu pada sembarang waktu $t$.
    
    \item \textbf{[C2 - Memahami]} Jelaskan secara fisis fenomena \textit{inductive kickback} saat sakelar sebuah sirkuit induktif DC dibuka mendadak. Mengapa polaritas tegangan induktor berbalik tanda saat sakelar dibuka?
    
    \item \textbf{[C3 - Menerapkan]} Sebuah kapasitor $C = 50\,\mu\text{F}$ yang memiliki tegangan awal $v_C(0) = 40\text{ V}$ dikosongkan melalui sebuah resistor $R = 200\text{ k}\Omega$. Tentukan persamaan tegangan kapasitor $v_C(t)$ dan hitung waktu yang dibutuhkan agar energi yang tersimpan di dalam kapasitor tersisa tinggal $25\%$ dari energi mula-mula.
    
    \item \textbf{[C4 - Menganalisis]} Tinjau rangkaian pengisian induktor dengan $V_0 = 100\text{ V}$, $R = 25\,\Omega$, dan $L = 5\text{ H}$. Tentukan persamaan laju disipasi daya sesaat pada resistor $p_R(t)$ dan laju penyerapan daya pada induktor $p_L(t)$. Pada waktu $t$ berapakah daya yang diserap oleh medan magnetik induktor mencapai nilai puncak maksimum?
    
    \item \textbf{[C5 - Mengevaluasi]} Sebuah transformator distribusi di kawasan industri memiliki konstanta waktu termal $\tau_{\text{th}} = 3\text{ jam}$. Trafo sedang beroperasi pada beban nominal dengan temperatur tunak $75^\circ\text{C}$ pada suhu ambien $30^\circ\text{C}$. Tiba-tiba terjadi lonjakan beban yang menyebabkan daya rugi-rugi meningkat menjadi 1,8 kali lipat dari kondisi normal. Jika batas aman temperatur puncak isolasi sebelum degradasi katastropik adalah $105^\circ\text{C}$, evaluasi apakah trafo tersebut aman beroperasi pada kondisi beban lebih ini selama 2 jam tanpa tindakan pemadaman darurat!
    
    \item \textbf{[C6 - Merancang]} Rancanglah sebuah skrip fungsi dalam bahasa \textbf{Julia} yang menerima parameter input matriks riwayat pembebanan dinamis per jam $P_{\text{load}}(t)$ selama 24 jam dan mengintegrasikan model termal diferensial transformator secara adaptif untuk menghitung total laju degradasi penuaan relatif ekuivalen ($L_{\text{loss}}$) menurut model IEEE Std C57.91.
\end{enumerate}

\section{Rangkuman \& Glosarium Istilah Teknis}

\subsection{Rangkuman Inti}
\begin{enumerate}[leftmargin=*]
    \item Seluruh fenomena transien orde 1 dalam rekayasa elektro dikendalikan oleh prinsip kontinuitas energi: $i_L(0^+) = i_L(0^-)$ pada induktor dan $v_C(0^+) = v_C(0^-)$ pada kapasitor.
    \item Respon waktu rangkaian orde 1 selalu mengikuti peluruhan eksponensial $e^{-t/\tau}$. Setelah selang waktu $t = 5\tau$, proses transien dianggap selesai ($>99{,}3\%$ kondisi tunak tercapai).
    \item Pemutusan arus mendadak pada beban induktif memicu lonjakan tegangan balik raksasa (\textit{inductive kickback}) yang berpotensi merusak sakelar. Mitigasi standar dilakukan dengan memasang dioda \textit{freewheeling} anti-paralel.
    \item Persamaan kesetimbangan energi kalor transformator memiliki struktur matematis yang identik dengan rangkaian RC paralel. Konstanta waktu termal $\tau_{\text{th}} = R_{\text{th}} C_{\text{th}}$ menentukan kecepatan respons suhu terhadap perubahan arus beban.
    \item Degradasi isolasi trafo mengikuti hukum Arrhenius, di mana setiap kenaikan suhu $6^\circ\text{C}$ di atas $110^\circ\text{C}$ menggandakan laju penuaan kertas selulosa.
\end{enumerate}

\subsection{Glosarium Istilah Teknis}
\begin{itemize}[leftmargin=*]
    \item \textbf{Inductive Kickback:} Lonjakan tegangan transien bernilai sangat tinggi yang dibangkitkan oleh induktor saat arus yang melaluinya diputus secara mendadak.
    \item \textbf{Freewheeling Diode (Flyback Diode):} Dioda penyearah yang dipasang secara paralel-kebalikan pada beban induktif untuk menyediakan jalur sirkulasi aman bagi arus pelepasan induktor.
    \item \textbf{Konstanta Waktu ($\tau$):} Waktu yang diperlukan oleh sistem orde 1 untuk menyelesaikan $63{,}2\%$ transisinya menuju kondisi tunak baru.
    \item \textbf{Hukum Pendinginan Newton:} Prinsip perpindahan kalor yang menyatakan bahwa laju kehilangan panas suatu benda berbanding lurus dengan selisih suhunya terhadap lingkungan sekitar.
    \item \textbf{Kapasitas Panas Efektif ($C_{\text{th}}$):} Jumlah energi kalor yang diperlukan untuk menaikkan temperatur massa total transformator sebesar satu derajat Celsius.
    \item \textbf{Hot-Spot Temperature:} Titik temperatur tertinggi pada belitan transformator daya yang menjadi parameter kritis penentu umur pakai isolasi.
\end{itemize}

\section{Referensi Standar Industri \& Buku Teks Acuan}
\begin{enumerate}[leftmargin=*]
    \item Alexander, C. K., \& Sadiku, M. N. O. (2021). \textit{Fundamentals of Electric Circuits} (7th ed.). New York: McGraw-Hill Education.
    \item IEEE Power and Energy Society. (2012). \textit{IEEE Std C57.91-2011: IEEE Guide for Loading Mineral-Oil-Immersed Transformers and Step-Voltage Regulators}. New York: IEEE.
    \item Boylestad, R. L. (2016). \textit{Introductory Circuit Analysis} (13th ed.). Boston: Pearson.
    \item Kreyszig, E. (2011). \textit{Advanced Engineering Mathematics} (10th ed.). Hoboken, NJ: John Wiley \& Sons.
    \item PT PLN (Persero). (2014). \textit{Standar Perusahaan Listrik Negara (SPLN) T5.003: Transformator Tenaga}. Jakarta: PT PLN (Persero).
    \item Rackauckas, C., \& Nie, Q. (2017). DifferentialEquations.jl -- A Performant and Feature-Rich Ecosystem for Solving Differential Equations in Julia. \textit{Journal of Open Research Software}, 5(1), 15.
\end{enumerate}

\end{document}
